\documentclass[english,hebrew]{article}
\usepackage[HE8,T1]{fontenc}
\usepackage[utf8]{inputenc}
\setcounter{secnumdepth}{-2}
\setlength{\parindent}{0bp}
\usepackage{calc}
\usepackage{amsmath}
\usepackage{amsthm}
\usepackage{amssymb}

\makeatletter
\theoremstyle{plain}
\newtheorem*{thm*}{\protect\theoremname}

\@ifundefined{date}{}{\date{}}
\usepackage{graphicx}
\usepackage{geometry}
\newcommand{\limi}{\underset{n\rightarrow\infty}{\lim}}
\newcommand{\limxz}{\underset{x\rightarrow x_{0}}{\lim}}
\newcommand{\limfxz}{\underset{x\rightarrow x_{0}}{\lim}f(x)}
\newcommand{\limfxm}{\underset{x\rightarrow x_{0}^{-}}{\lim}f(x)}
\newcommand{\limfxp}{\underset{x\rightarrow x_{0}^{+}}{\lim}f(x)}
\newcommand{\limxm}{\underset{x\rightarrow x_{0}^{-}}{\lim}}
\newcommand{\limxp}{\underset{x\rightarrow x_{0}^{+}}{\lim}}
\newcommand{\sqrm}{M_{m\times m}(\mathbb{F})}
\newcommand{\seqa}{(a_{n})_{n=1}^{\infty}}
\newcommand{\seqx}{(x_{n})_{n=1}^{\infty}}
\newcommand{\funcdr}{f:D\rightarrow\mathbb{R}}
\newcommand{\der}{\limxz\frac{f(x)-f(x_{0})}{x-x_{0}}}
\usepackage[all]{pdfprivacy}

\makeatother

\usepackage{babel}
\addto\captionsenglish{\renewcommand{\theoremname}{Theorem}}
\addto\captionshebrew{\renewcommand{\theoremname}{משפט}}
\providecommand{\theoremname}{משפט}

\begin{document}
\newgeometry{top=1cm, bottom=2cm, left=1cm, right=1cm}
\title{אלגברה ליניארית {\beginL 2\endL} - הרצאה {\beginL 2\endL}}
\author{נעם לב}
\date{\L{$14.10.2026$}}

\maketitle
\noindent\fbox{\begin{minipage}[t]{1\columnwidth - 2\fboxsep - 2\fboxrule}%
\begin{center}
האוניברסיטה העברית \textbar{} אלגברה ליניארית ){\beginL 2\endL}( -
\L{$80135$} \textbar{} סמסטר{\small{} }{\beginL 2027\endL}{\small{} }א'
)תשפ\textquotedblleft ז א'( \textbar{} מרצה: פרופסור יבגני סטרחוב
\textbar{} מתרגל: מר לזר גוטרמן
\par\end{center}%
\end{minipage}}
\begin{center}
הסיכום הוא לא רשמי ולא עבר אישור של הסגל, על כן אין להשתמש בו )ובין
השאר במספור שלו ובשמות הטענות שבו( כמקור רשמי. אם ישנן טעויות, הן
שלי ולא של הסגל.
\par\end{center}

\medskip{}

\tableofcontents{}

\newpage{}

\part{ערכים עצמיים}

\rule[0.5ex]{1\columnwidth}{1pt}

\section{תזכורת}

\subsection{הגדרה - ערך עצמי של אופרטור \L{$T$} - \L{$\text{Eigenvalue of }T$}}

יהי \L{$V$} מ\textquotedblright ו מעל \L{$\mathbb{F}$}.

עבור אופרטור \L{$T\in\mathcal{L}(V)$}, אם קיים וקטור \L{$0_{V}\ne v\in V$}
כך שמתקיים \L{%
\[
T(v)=\lambda v
\]
}%
כאשר \L{$\lambda\in\mathbb{F}$}, נאמר כי \L{$\lambda$} הוא ערך עצמי
)\L{$\text{Eigenvalue}$}( של \L{$T$}.

\subsection{הגדרה - וקטור עצמי}

לכל \L{$v$} כזה נקרא וקטור עצמי של \L{$T$} השייך ל-\L{$\lambda$}. 

\rule[0.5ex]{1\columnwidth}{1pt}

\section{משפט - אפיון ערכים עצמיים באמצעות דטרמיננטה}

יהי \L{$V$} מ\textquotedblright ו מעל \L{$\mathbb{F}$} כך ש- \L{%
\[
\dim V=n\geq1
\]
}%
\L{%
\[
T\in\mathcal{L}(V)
\]
}%
\L{%
\[
\mathcal{B}_{V}=(v_{1},\dots,v_{n})
\]
}%

)\L{$\mathcal{B}_{V}$} בסיס של \L{$V$}(. 

אז התנאים הבאים שקולים: 
\begin{enumerate}
\item \L{$\lambda\in\mathbb{F}$} הוא ערך עצמי של \L{$T$}. 
\item מתקיים \L{%
\[
\det(M(T,\mathcal{B}_{V},\mathcal{B}_{V})-\lambda I_{n})=0_{\mathbb{F}}
\]
}%
\end{enumerate}
\begin{proof}
1. \L{$\lambda$} - ערך עצמי של \L{$T$} \L{$\iff$} קיים \L{$0_{V}\ne v\in V$}
כך ש-\L{$T(v)=\lambda v$}.

{\beginL 2\endL}. \L{$\mathcal{B}_{V}$} בסיס של \L{$V$} לכן ניתן
להציג את \L{$v$} כצ\textquotedblright ל שלו\L{%
\[
v=b_{1}v_{1}+\dots+b_{n}v_{n}
\]
}%

מכאן סעיף {\beginL 1\endL} מתקיים אם\textquotedblright ם קיימים \L{%
\[
b_{1},\dots,b_{n}\in\mathbb{F}:\exists1\le i\le n\text{ s.t. }b_{i}\ne0
\]
}%
 כך ש-\L{%
\[
T\left(b_{1}v_{1}+\dots+b_{n}v_{n}\right)=\lambda\left(b_{1}v_{1}+\dots+b_{n}v_{n}\right)=\lambda b_{1}v_{1}+\dots+\lambda b_{n}v_{n}
\]
}%
{\beginL 3\endL}. נסתכל על \L{$M(T,\mathcal{B}_{V},\mathcal{B}_{V})$}\L{%
\[
M(T,\mathcal{B}_{V},\mathcal{B}_{V})=\begin{pmatrix}a_{1,1} & \cdots & a_{1,n}\\
\vdots &  & \vdots\\
a_{n,1} & \cdots & a_{n,n}
\end{pmatrix}
\]
}%

{\beginL 4\endL}. לפי לינאריות ו-{\beginL 3\endL} מתקיים \L{%
\[
\begin{cases}
T\left(b_{1}v_{1}+\dots+b_{n}v_{n}\right) & \underset{\text{linearity}}{=}b_{1}T(v_{1})+\dots+b_{n}T(v_{n})\\
 & \underset{3}{=}b_{1}(a_{1,1}v_{1}+\dots+a_{n,1}v_{n})+\dots+b_{n}(a_{1,n}v_{1}+\dots+a_{n,n}v_{n})\\
 & =\left(b_{1}a_{1,1}+\dots+b_{n}a_{1,n}\right)v_{1}+\dots+\left(b_{1}a_{n,1}+\dots+b_{n}a_{n,n}\right)v_{n}
\end{cases}
\]
}%
לכן, לפי המשפט שלמדנו בלינארית {\beginL 1\endL} על הצגה יחידה של וקטור
לפי בסיס, נקבל ש-\L{$T(v)=\lambda v$} אם\textquotedblright ם \L{%
\[
\begin{cases}
b_{1}a_{1,1}+\dots+b_{n}a_{1,n}=\lambda b_{1}\\
\quad\vdots\\
b_{1}a_{n,1}+\dots+b_{n}a_{n,n}=\lambda b_{n}
\end{cases}
\]
}%
\L{%
\[
\iff\begin{pmatrix}a_{1,1} & \cdots & a_{1,n}\\
\vdots &  & \vdots\\
a_{n,1} & \cdots & a_{n,n}
\end{pmatrix}\begin{pmatrix}b_{1}\\
\vdots\\
b_{n}
\end{pmatrix}=\lambda\begin{pmatrix}b_{1}\\
\vdots\\
b_{n}
\end{pmatrix}
\]
}%
\L{%
\[
\iff M(T,\mathcal{B}_{V},\mathcal{B}_{V})\cdot\begin{pmatrix}b_{1}\\
\vdots\\
b_{n}
\end{pmatrix}=\lambda\begin{pmatrix}b_{1}\\
\vdots\\
b_{n}
\end{pmatrix}
\]
}%
\L{%
\[
\iff\left(M(T,\mathcal{B}_{V},\mathcal{B}_{V})-\lambda I_{n}\right)\cdot\begin{pmatrix}b_{1}\\
\vdots\\
b_{n}
\end{pmatrix}=\begin{pmatrix}0\\
\vdots\\
0
\end{pmatrix}
\]
}%

לכן \L{$\lambda$} ערך עצמי של \L{$T$} אם\textquotedblright ם לממ\textquotedblright ל
ההומוגנית הזאת יש פתרון לא טריוויאלי )כי \L{$\exists1\le i\le n\text{ s.t }b_{i}\ne0$}(.

לממ\textquotedblright ל זה יש פתרון לא טריוויאלי אם\textquotedblright ם
המטריצה \L{$M(T,\mathcal{B}_{V},\mathcal{B}_{V})-\lambda I_{n}$}
לא הפיכה, שזה מתקיים אם\textquotedblright ם \L{%
\[
\det\left(M(T,\mathcal{B}_{V},\mathcal{B}_{V})-\lambda I_{n}\right)=0_{\mathbb{F}}
\]
}%

כנדרש.
\end{proof}
הערות:

- נוח לעשות כאן את שני הכיוונים יחד, כי כל המעברים הם שקילויות.

- המעבר האחרון הוא שקילות שלמדנו בלינארית {\beginL 1\endL}.

- הדרישה \L{$n\ge1$} היא כי אם \L{$\dim V=0$} אז \L{$V=\{0_{V}\}$},
ו-\L{$0_{V}$} אינו וקטור עצמי. כעיקרון, הטענה מתקיימת כאן באופן ריק,
אבל זה לא מעניין. 

\section{הגדרה - הפולינום האופייני}

\L{%
\[
p(\lambda)\overset{\text{def}}{=}\det(\lambda I_{n}-M(T,\mathcal{B}_{V},\mathcal{B}_{V}))
\]
}%

1. \L{$p(\lambda)$} אינו תלוי בבחירה של \L{$\mathcal{B}_{V}$}.

2. \L{$p(\lambda)$} הוא פולינום עם מקדמים ב-\L{$\mathbb{F}$}.

3. \L{$\deg(p(\lambda))\leq n$}.

)בחישוב הדטרמיננטה החזקה הגבוהה ביותר שאפשר לקבל היא \L{$\lambda^{n}$},
ולכן המעלה היא לכל היותר \L{$n$}.

למעשה, נראה בהמשך ש-\L{$\deg(p(\lambda))=n$} בדיוק, והמקדם המוביל
הוא \L{$1$}.(

\subsection{למה הפולינום האופייני לא תלוי בבחירה של הבסיס?}

אם \L{$\mathcal{\mathcal{C}_{V}}$} בסיס אחר של \L{$V$}, אז\L{%
\[
M(T,\mathcal{\mathcal{C}_{V}},\mathcal{\mathcal{C}_{V}})=P^{-1}M(T,\mathcal{B}_{V},\mathcal{B}_{V})P
\]
}%

עבור מטריצת המעבר ההפיכה \L{$P=[\text{Id}_{V}]_{\mathcal{B}_{V}}^{\mathcal{\mathcal{C}_{V}}}$}.
לכן,\L{%
\[
\lambda I_{n}-M(T,\mathcal{\mathcal{C}_{V}},\mathcal{\mathcal{C}_{V}})=P^{-1}(\lambda I_{n}-M(T,\mathcal{B}_{V},\mathcal{B}_{V}))P
\]
}%

ולמטריצות דומות יש אותה דטרמיננטה.\\

לכן מותר לדבר על \textquotedblright הפולינום האופייני של \L{$T$}\textquotedblright . 

\subsection{אלגוריתם למציאת ערכים עצמיים}

שלב {\beginL 1\endL}: לחשב מטריצה מייצגת בבסיס כלשהו.

שלב {\beginL 2\endL}: למצוא את הפולינום האופייני.

שלב {\beginL 3\endL}: למצוא את השורשים של \L{$p(\lambda)$} בשדה \L{$\mathbb{F}$}.
לפי המשפט הקודם, אלה בדיוק הערכים העצמיים של \L{$T$}.

\subsection{\L{$\det(\lambda I_{n}-M)$} מול \L{$\det(M-\lambda I_{n})$}}

במשפט כתבנו \L{$\det(M-\lambda I_{n})$} ואילו כשחישבנו כתבנו \L{$\det(\lambda I_{n}-M)$},
כיוון שבצורה הראשונה יותר נוח להוכיח, ובצורה השנייה יותר נוח לחשב.\L{%
\[
\lambda I_{n}-M=-(M-\lambda I_{n})
\]
}%

וזה שקול להכפלת כל שורה של \L{$(M-\lambda I_{n})$} במינוס אחת.

הדטרמיננטה היא פונ' נפח ולכן לינארית בכל שורה, מכאן הכפלת כל אחת מהשורות
ב-\L{$(-1)$} נותנת \L{%
\[
\det(\lambda I_{n}-M)=(-1)^{n}\det(M-\lambda I_{n})
\]
}%
מאחר ש-\L{$(-1)^{n}\neq0$}, הביטוי \L{$\det(M-\lambda I_{n})$} מתאפס
אם\textquotedblright ם הביטוי \L{$\det(\lambda I_{n}-M)$} מתאפס.

אנחנו משווים לאפס, ולכן החישוב של שניהם שקול.

\subsection{דוגמה - מציאת ערכים עצמיים}

\L{%
\[
V=\mathbb{C}^{2}=\{(z_{1},z_{2})\mid z_{1}\in\mathbb{C},z_{2}\in\mathbb{C}\}
\]
}%
\L{%
\[
T((z_{1},z_{2}))\overset{\text{def}}{=}(-z_{2},z_{1})
\]
}%
\L{%
\[
T\in\mathcal{L}(\mathbb{C}^{2})
\]
}%
חשבו את כל הערכים העצמיים של \L{$T$}.

\subsubsection{פתרון}

{\beginL 1\endL}. נבחר את הבסיס הסטנדרטי\L{%
\[
\mathcal{B}_{V}=\left(\underset{v_{1}}{(1,0)},\underset{v_{2}}{(0,1)}\right)
\]
}%

אז \L{%
\begin{align*}
T(v_{1}) & =T((1,0))=(0,1)=0\cdot v_{1}+1\cdot v_{2}\\
T(v_{2}) & =T((0,1))=(-1,0)=-1\cdot v_{1}+0\cdot v_{2}
\end{align*}
}%
לכן\L{%
\[
M(T,\mathcal{B}_{V},\mathcal{B}_{V})=\begin{pmatrix}0 & -1\\
1 & 0
\end{pmatrix}
\]
}%
)תזכורת - העמודה הראשונה היא המקדמים של \L{$T(v_{1})$} והעמודה השנייה
היא המקדמים של \L{$T(v_{2})$}(\\

2. \L{%
\[
\det\left[\lambda\begin{pmatrix}1 & 0\\
0 & 1
\end{pmatrix}-\begin{pmatrix}0 & -1\\
1 & 0
\end{pmatrix}\right]=\det(\lambda I-M(T,\mathcal{B}_{V},\mathcal{B}_{V}))
\]
}%
\L{%
\[
=\det\begin{pmatrix}\lambda & 1\\
-1 & \lambda
\end{pmatrix}=\lambda\cdot\lambda-1\cdot\left(-1\right)=\lambda^{2}+1=0
\]
}%
\L{%
\[
\lambda_{1}=i,\lambda_{2}=-i
\]
}%
הם הערכים העצמיים של \L{$T$}.

\subsection{תלות בשדה}

אם במקום \L{$\mathbb{C}$} בוחרים את \L{$\mathbb{R}$}, כלומר \L{$V=\mathbb{R}^{2}$}
ואותה העתקה, אז ל-\L{$\lambda^{2}+1$} אין שורשים ממשיים, ולכן אין
ערכים עצמיים בכלל - השדה שבוחרים משנה.

\subsection{המחשה על למה \L{$\deg(p(\lambda))=n$} והמקדם המוביל הוא \L{$a_{n}=1$}}

נמחיש בעזרת \L{$n=3$}. אז המטריצה שעליה מפעילים דטרמיננטה בפולינום
האופייני היא מהצורה\L{%
\[
\begin{pmatrix}\lambda-a_{1,1} & -a_{1,2} & -a_{1,3}\\
-a_{2,1} & \lambda-a_{2,2} & -a_{2,3}\\
-a_{3,1} & -a_{3,2} & \lambda-a_{3,3}
\end{pmatrix}
\]
}%
ומתקיים

\L{%
\[
\det\begin{pmatrix}\lambda-a_{1,1} & -a_{1,2} & -a_{1,3}\\
-a_{2,1} & \lambda-a_{2,2} & -a_{2,3}\\
-a_{3,1} & -a_{3,2} & \lambda-a_{3,3}
\end{pmatrix}=
\]
}%
\L{%
\[
\det\left(\left(\lambda-a_{1,1}\right)\cdot\begin{pmatrix}\lambda-a_{2,2} & -a_{2,3}\\
-a_{3,2} & \lambda-a_{3,3}
\end{pmatrix}-\left(-a_{1,2}\right)\cdot\begin{pmatrix}-a_{2,1} & -a_{2,3}\\
-a_{3,1} & \lambda-a_{3,3}
\end{pmatrix}+\left(-a_{1,3}\right)\cdot\begin{pmatrix}-a_{2,1} & \lambda-a_{2,2}\\
-a_{3,1} & -a_{3,2}
\end{pmatrix}\right)
\]
}%

המחובר הראשון נותן פולינום שמתחיל ב-\L{$\lambda^{3}$} והמחוברים השני
והשלישי נותנים פולינומים ממעלה לכל היותר {\beginL 1\endL}.

לכן, החזקה המקסימלית היא \L{$\lambda^{3}$}, והמקדם שלה הוא \L{$1$}.\\

באופן כללי: בפיתוח המלא של הדטרמיננטה, מכפלת איברי האלכסון \L{$(\lambda-a_{1,1})\cdots(\lambda-a_{n,n})$}
היא היחידה שמגיעה ל-\L{$\lambda^{n}$}.

כל מכפלה אחרת מדלגת על לפחות שני איברי אלכסון, ולכן היא ממעלה \L{$\leq n-2$}.

לכן \L{%
\[
p(\lambda)=\lambda^{n}+\dots
\]
}%
 מעל כל שדה.

דהיינו, המעלה של הפולינום האופייני תמיד שווה למימד, מעל כל שדה.

החזקה הגבוהה ביותר מתקבלת תמיד מהאלכסון של המטריצה, ושאר המחוברים
נותנים מעלה נמוכה יותר.

\part{המשפט היסודי של האלגברה - \L{$\text{The Fundamental Theorem of Algebra}$}}
\begin{thm*}
לכל פולינום )שאינו קבוע( עם מקדמים מרוכבים יש לפחות שורש מרוכב אחד,
ולכן ניתן לפרק אותו לגורמים לינאריים כך שמספר השורשים הכולל )כולל
חזרות( שווה בדיוק למעלת הפולינום.
\end{thm*}
זו תוצאה שלא מוכיחים בקורס זה, רק משתמשים בה.

זו טענה מאנליזה, ומוכיחים אותה למשל בקורס פונקציות מרוכבות.

יהי \L{%
\[
P_{n}(z)=a_{n}z^{n}+a_{n-1}z^{n-1}+\dots+a_{1}z+a_{0}
\]
}%
\L{%
\[
\text{s.t. }a_{n},\dots,a_{1},a_{0}\in\mathbb{C},\quad a_{n}\neq0,\quad n\geq1
\]
}%
אז ניתן לכתוב: \L{%
\[
P_{n}(z)=a_{n}(z-\lambda_{1})^{m_{1}}(z-\lambda_{2})^{m_{2}}\dots(z-\lambda_{k})^{m_{k}}
\]
}%
עם \L{$1\le k\le n$} ומתקיים \L{%
\[
m_{1},\dots,m_{k}\in\mathbb{N},\quad m_{1}+\dots+m_{k}=n,\quad\forall1\le i\le k:m_{i}\ge1,\quad\lambda_{1},\dots,\lambda_{k}\in\mathbb{C}
\]
}%

\begin{itemize}
\item \L{$\lambda_{1},\dots,\lambda_{k}$} שונים זה מזה.
\item \L{$m_{i}$} נקרא הריבוי של השורש \L{$\lambda_{i}$}.
\item מאחר ש-\L{$k\geq1$}, לפולינום יש לפחות שורש אחד.
\item הפולינום האופייני הוא מקרה פרטי, עם \L{$a_{n}=1$}.\\
עקרונית, אפשר להוכיח את קיום הערך העצמי גם בלי דטרמיננטות )כך עושה
\L{$\text{Axler}$}(, אבל אין דרך לברוח מהמשפט היסודי של האלגברה:
תמיד צריך תוצאה כלשהי מאנליזה.
\end{itemize}

\section{דוגמה}

\L{%
\[
p(z)=z^{4}+1
\]
}%
השורשים מקיימים \L{%
\[
z^{4}=-1=e^{i\left(\pi+2\pi k\right)}:k\in\mathbb{Z}
\]
}%
\L{%
\[
z=e^{\frac{i\left(\pi+2\pi k\right)}{4}}:k\in\mathbb{Z}
\]
}%
 על מעגל היחידה יש ארבע נקודות כאלה: \L{%
\[
e^{i\frac{\pi}{4}},\quad e^{i\frac{3\pi}{4}},\quad e^{i\frac{5\pi}{4}},\quad e^{i\frac{7\pi}{4}}
\]
}%

לכן\L{%
\[
p(z)=(z-e^{i\frac{\pi}{4}})(z-e^{i\frac{3\pi}{4}})(z-e^{i\frac{5\pi}{4}})(z-e^{i\frac{7\pi}{4}})
\]
}%

הריבוי של כל שורש כאן הוא {\beginL 1\endL}.

\section{משפט - לכל אופרטור על מ\textquotedblright ו נ\textquotedblright ס
מעל \L{$\mathbb{C}$} יש ערך עצמי}

)יבגני אמר שהמשפט הזה כנראה יהיה במבחן, ושזו \textquotedblright תופעה
של טבע\textquotedblright , כמו חוק ניוטון: אחד הדברים הבסיסיים(.\\

יהי \L{$V$} מ\textquotedblright ו נ\textquotedblright ס מעל \L{$\mathbb{C}$}
כך ש- \L{%
\[
\dim V=n\geq1
\]
}%
\L{%
\[
T\in\mathcal{L}(V)
\]
}%
אז קיים ערך עצמי של \L{$T$} )לפחות אחד(.
\begin{proof}
1. נבחר בסיס \L{$\mathcal{B}_{V}=(v_{1},\dots,v_{n})$} של \L{$V$}.

2. לפי המשפט הקודם, \L{$\lambda$} ערך עצמי של \L{$T$} \L{$\iff$}\L{%
\[
p(\lambda)=\det(\lambda I-M(T,\mathcal{B}_{V},\mathcal{B}_{V}))=0
\]
}%
3. \L{%
\[
p(\lambda)=\lambda^{n}+b_{n-1}\lambda^{n-1}+\dots+b_{1}\lambda+b_{0}
\]
}%
\L{%
\[
b_{0},b_{1},\dots,b_{n-1}\in\mathbb{C}
\]
}%
\L{%
\[
\deg(p(\lambda))=n\geq1
\]
}%

4. לפי המשפט היסודי של האלגברה, לפולינום ממעלה \L{$n\geq1$} מעל \L{$\mathbb{C}$}
יש שורש \L{$\lambda_{1}\in\mathbb{C}$}. לפי 2, \L{$\lambda_{1}$}
הוא ערך עצמי של \L{$T$}. 
\end{proof}

\subsection{אנטי-דוגמה - אופרטור ללא ערכים עצמיים )מימד אינסופי(}

ראינו שהשדה \L{$\mathbb{C}$} חשוב.

האם גם המימד הסופי הכרחי?

נראה העתקה לינארית במרחב וקטורי מעל \L{$\mathbb{C}$}, ממימד אינסופי,
בלי ערכים עצמיים.\\

\L{$V=\mathcal{P}(\mathbb{C})$} - כל הפולינומים מעל \L{$\mathbb{C}$}.
\L{%
\[
(T(p))(z)=z\,p(z)
\]
}%
\L{%
\[
T\in\mathcal{L}(\mathcal{P}(\mathbb{C}))
\]
}%
האופרטור הזה מכפיל את הפולינום במשתנה.

למשל \L{$T\left(z^{2}+1\right)=z^{3}+z$}.

\L{%
\[
(T(p+q))(z)=z(p(z)+q(z))=z\,p(z)+z\,q(z)=(T(p))(z)+(T(q))(z)
\]
}%
\L{%
\[
(T(cp))(z)=z\,c\,p(z)=c\,(T(p))(z)
\]
}%
\textbf{טענה:} אין ערכים עצמיים של \L{$T$}.
\begin{proof}
נניח בשלילה כי \L{$\lambda$} ערך עצמי, כלומר קיים \L{$p\neq0$} )פולינום
שאינו פולינום האפס( כך ש- \L{%
\[
(T(p))(z)=\lambda p(z)
\]
}%
\L{%
\[
\Rightarrow z\,p(z)=\lambda p(z)
\]
}%
\L{%
\[
\Rightarrow(z-\lambda)p(z)=0
\]
}%

היות ו-\L{$p$} הוא לא פולינום האפס,\L{%
\[
\deg((z-\lambda)p(z))=1+\deg p\geq1
\]
}%

לכן \L{$(z-\lambda)p(z)$} הוא לא פולינום האפס כלומר \L{%
\[
(z-\lambda)p(z)\ne0
\]
}%

והגענו לסתירה.
\end{proof}

\section{משפט - וקטורים עצמיים השייכים לערכים עצמיים שונים הם בת\textquotedblright ל}

)יבגני אמר שגם המשפט הזה כנראה יופיע במבחן, הוא מאוד בסיסי.

הוא חשוב כי הוא מאפשר, בתנאים מסוימים, לבנות בסיסים של וקטורים עצמיים,
ובבסיסים כאלה המטריצה של ההעתקה מאוד פשוטה. שימו לב: במשפט הזה מדובר
על כל שדה ועל כל מרחב, כלומר אין הגבלות(\\

יהי \L{$V$} מ\textquotedblright ו מעל \L{$\mathbb{F}$} ו-\L{$T\in\mathcal{L}(V)$}
כך ש-\L{%
\begin{align*}
T(v_{1}) & =\lambda_{1}v_{1}\\
 & \vdots\\
T(v_{m}) & =\lambda_{m}v_{m}
\end{align*}
}%

כאשר \L{$\lambda_{1},\dots,\lambda_{m}$} שונים זה מזה וגם\L{%
\[
\forall1\le i\le m:\lambda_{i}\in\mathbb{F}\wedge0_{V}\ne v_{i}\in V
\]
}%
כלומר לכל \L{$i$}, \L{$v_{i}$} הוא וקטור עצמי השייך ל-\L{$\lambda_{i}$}.\\

אז \L{$(v_{1},\dots,v_{m})$} בלתי תלויה לינארית.
\begin{proof}
באינדוקציה על \L{$m$}.

1. \L{$m=1$}: הסדרה \L{$(v_{1})$} היא בלתי תלויה לינארית, כי \L{$v_{1}\neq0_{V}$}.

2. נניח כי המשפט מתקיים עבור \L{$m=k-1$}. נוכיח כי הוא מתקיים גם
עבור \L{$m=k$}.

3. נניח בשלילה כי \L{$(v_{1},\dots,v_{m})$} תלויה לינארית. אז קיימים
\L{$d_{1},\dots,d_{m}\in\mathbb{F}$} כך ש- \L{%
\[
d_{1}v_{1}+\dots+d_{m}v_{m}=0_{V}
\]
}%
ולא כולם שווים ל-\L{$0_{\mathbb{F}}$}.

4. נפעיל אופרטור חדש, \L{%
\[
\left(T-\lambda_{m}I\right)\left(v\right)=T(v)-\lambda_{m}I(v)=T(v)-\lambda_{m}v
\]
}%
על שני האגפים: \L{%
\begin{align*}
0_{V} & =(T-\lambda_{m}I)(d_{1}v_{1}+\dots+d_{m-1}v_{m-1}+d_{m}v_{m})\\
 & =d_{1}\cdot(T-\lambda_{m}I)(v_{1})+\dots+d_{m}\cdot(T-\lambda_{m}I)(v_{m})\\
 & =d_{1}\cdot\left(T(v_{1})-\lambda_{m}v_{1}\right)+\dots+d_{m}\cdot\left(T(v_{m})-\lambda_{m}v_{m}\right)\\
 & =d_{1}\cdot\left(\lambda_{1}v_{1}-\lambda_{m}v_{1}\right)+\dots+d_{m}\cdot\left(\lambda_{m}v_{m}-\lambda_{m}v_{m}\right)\\
 & =d_{1}\underset{\neq0_{\mathbb{F}}}{\underbrace{(\lambda_{1}-\lambda_{m})}}v_{1}+\dots+d_{m-1}\underset{\neq0_{\mathbb{F}}}{\underbrace{(\lambda_{m-1}-\lambda_{m})}}v_{m-1}+\underset{=0_{V}}{\underbrace{d_{m}(\lambda_{m}-\lambda_{m})v_{m}}}
\end{align*}
}%
)\L{$\lambda_{1},\dots,\lambda_{m}$} שונים זה מזה, ולכן ההפרשים שונים
מ-\L{$0_{\mathbb{F}}$}(.

5. קיבלנו צ\textquotedblright ל של \L{$(v_{1},\dots,v_{m-1})$} ששווה
ל-\L{$0_{V}$}.

לפי הנחת האינדוקציה \L{$(v_{1},\dots,v_{m-1})$} בלתי תלויה לינארית,
ולכן \L{%
\[
\forall1\le i\le m-1:d_{i}(\lambda_{i}-\lambda_{m})=0_{\mathbb{F}}
\]
}%
 ומאחר ש- \L{$\lambda_{i}-\lambda_{m}\neq0_{\mathbb{F}}$},\L{%
\[
d_{1}=\dots=d_{m-1}=0_{\mathbb{F}}
\]
}%

6. נציב בשוויון מסעיף 3 ונקבל \L{$d_{m}v_{m}=0_{V}$}.

אבל \L{$v_{m}\neq0_{V}$}, לכן \L{$d_{m}=0_{\mathbb{F}}$}.

כלומר, כל המקדמים שווים ל- \L{$0_{\mathbb{F}}$}, בסתירה לסעיף 3. 
\end{proof}
מסקנה )לקראת ההמשך(: אם \L{$\dim V=n$}, אז בסדרה בת\textquotedblright ל
יש לכל היותר \L{$n$} וקטורים, ולכן ל-\L{$T$} יש לכל היותר \L{$n$}
ערכים עצמיים שונים.

אם יש ל-\L{$T$} בדיוק \L{$n$} ערכים עצמיים שונים \L{$\lambda_{1},\dots,\lambda_{n}$},
אז וקטורים עצמיים מתאימים \L{$\mathcal{B}_{V}=(v_{1},\dots,v_{n})$}
מהווים בסיס של \L{$V$} )סדרה בת\textquotedblright ל באורך \L{$\dim V$}(,
ובבסיס הזה \L{$T(v_{j})=\lambda_{j}v_{j}$}, ולכן \L{$M(T,\mathcal{B}_{V},\mathcal{B}_{V})$}
אלכסונית עם \L{$\lambda_{1},\dots,\lambda_{n}$} על האלכסון. זו \textquotedblright המטריצה
המאוד פשוטה\textquotedblright{} שיבגני דיבר עליה.
\end{document}
